\partstart{Đặt vấn đề và phân tích yêu cầu}
\subsection{Bối cảnh và động cơ nghiên cứu}
Trong chuỗi cung ứng, một đơn hàng đi qua nhiều bước từ tiếp nhận, chuẩn bị
hàng, lựa chọn phương thức vận chuyển đến giao cho khách. Nếu chỉ phát hiện
chậm giao khi đã vượt thời hạn cam kết, cơ hội điều phối sớm bị thu hẹp.
Bài toán của đề tài 16 vì vậy tập trung vào cảnh báo nguy cơ: sử dụng các
thuộc tính có thể biết khi tạo đơn để nhận diện những đơn cần được theo dõi.
Kết quả phân loại có thể là đầu vào cho nhân viên vận hành xem xét, thay vì
tự động thay đổi cam kết giao hàng.

Một cảnh báo đúng có thể hỗ trợ ưu tiên kiểm tra tồn kho hoặc rà soát kế
hoạch vận chuyển. Ngược lại, cảnh báo sai gây thêm khối lượng xử lý, còn bỏ
sót khiến đơn có vấn đề không được chú ý. Báo cáo không gán giá trị tiền cho
những tác động này vì dữ liệu hiện có không chứa chi phí can thiệp. Sự hữu
ích của cảnh báo phải được đánh giá cùng khả năng xử lý của đơn vị vận hành,
không thể suy ra chỉ từ độ chính xác của mô hình.

Đề tài phù hợp với dữ liệu bảng có cả số và phân loại. K-NN được chọn để
khảo sát việc dự đoán từ những đơn có đặc trưng tương tự; cây quyết định
được chọn để khảo sát quy tắc chia nhóm. Hai thuật toán có cách biểu diễn
ranh giới khác nhau, tạo cơ sở so sánh giữa học dựa trên khoảng cách và học
dựa trên các điều kiện phân nhánh. Đây cũng là hai phương pháp thuộc phạm
vi thuật toán học phần quy định \cite{fit}.

\subsection{Mục tiêu, thời điểm dự đoán và phạm vi}
Biến mục tiêu là \code{Late_delivery_risk}, nhận giá trị 1 nếu nguồn gán
nguy cơ giao trễ, và 0 cho phần còn lại. Bài toán được phát biểu dưới dạng
học có giám sát: từ vector đặc trưng $x_i$ của một dòng hàng, ước lượng nhãn
$\hat y_i\in\{0,1\}$. Mỗi dòng còn gắn với một mã đơn $g_i$; các dòng chung
mã đơn có thể dùng chung phương thức vận chuyển và thông tin khách hàng.

Thời điểm dự đoán giả định là sau khi đơn được tạo và phương thức giao đã
được chọn, trước ngày vận chuyển và kết quả giao thực tế. Giả định này quyết
định các trường được phép làm đầu vào. Nếu doanh nghiệp muốn dự đoán trước
cả bước lựa chọn phương thức giao, danh sách đặc trưng hiện tại không còn
phù hợp và phải thiết kế lại bài toán. Tính sẵn có của từng trường cần được
chủ dữ liệu xác nhận trước ứng dụng.

Mục tiêu thực nghiệm gồm xây dựng quy trình chống rò rỉ, so sánh baseline
với K-NN và cây quyết định, đối chiếu hai luồng đặc trưng, đồng thời phân
tích lỗi theo phương thức giao hàng. Phạm vi hiện tại là phân loại nhị phân
trên dữ liệu lịch sử, không dự báo số ngày trễ, không tối ưu tuyến đường và
không chứng minh hiệu quả của một biện pháp can thiệp. Đánh giá theo đơn
được bổ sung để kiểm tra ý nghĩa của đơn vị quan sát.

\subsection{Nguồn và cấu trúc DataCo}
Theo trang nguồn Mendeley, DataCo bao gồm dữ liệu chuỗi cung ứng dạng bảng
và dữ liệu truy cập; tệp mô tả biến đi kèm giúp tra cứu ý nghĩa cột
\cite{dataco}. Đề tài sử dụng bản dữ liệu bảng ở liên kết Kaggle được lưu
trong manifest, không dùng clickstream và không ghép thêm nguồn ngoài
\cite{kaggle}. Việc có nhiều loại tệp ở nguồn gốc không có nghĩa tất cả đều
được đưa vào thực nghiệm.

Kiểm toán toàn bộ CSV ghi nhận 180.519 dòng, 53 cột và 65.752 mã đơn phân
biệt. Có 45.902 đơn chứa nhiều dòng; số dòng lặp mã đơn là 114.767. Đây là
cấu trúc đơn--mặt hàng, không phải 114.767 bản ghi lỗi. Không có dòng trùng
hoàn toàn và không có đơn chứa nhãn mâu thuẫn giữa các dòng, theo
\code{quality.json}. Hai kiểm tra này có ý nghĩa khác nhau: trùng toàn bộ
kiểm tra bản ghi, còn nhất quán nhãn kiểm tra khả năng tổng hợp theo đơn.

\begin{table}[!htbp]\centering
\caption{Các nhóm thông tin trong bộ dữ liệu và cách sử dụng.}\label{tab:data-groups}
\begin{tabular}{p{5.1cm}p{8.2cm}}\toprule
Nhóm trường & Ví dụ và vai trò\\\midrule
Nhãn và kết quả & Late delivery risk; trạng thái giao; số ngày giao thực tế. Chỉ nhãn dùng làm đích, trường kết quả bị loại khỏi đầu vào.\\
Đơn và sản phẩm & Số lượng, giá sản phẩm, tỷ lệ chiết khấu, nhóm sản phẩm. Một đơn có thể chứa nhiều mặt hàng.\\
Vận chuyển và thị trường & Phương thức giao, số ngày kế hoạch, vùng và quốc gia nhận đơn.\\
Định danh và cá nhân & Mã đơn, mã khách, tên, địa chỉ, email. Mã đơn chỉ dùng gom nhóm, thông tin cá nhân không dùng để dự đoán.\\\bottomrule
\end{tabular}

\end{table}
Bảng~\ref{tab:data-groups} cho thấy cùng một CSV có thể chứa cả thông tin
đầu vào hợp lệ lẫn thông tin hình thành sau kết quả. Vì vậy, đưa toàn bộ cột
vào mô hình sẽ làm thay đổi ý nghĩa dự báo sớm, ngay cả khi chương trình vẫn
chạy và metric tăng. Báo cáo dùng danh sách cho phép để kiểm soát vấn đề này.

\subsection{Phân phối nhãn và giới hạn ngữ nghĩa}
Toàn bộ nguồn có 98.977 dòng dương và 81.542 dòng âm, tỷ lệ dương khoảng
54,83\%. Hai lớp không mất cân bằng cực đoan, nhưng lớp dương là đa số. Điều
này khiến chiến lược luôn dự đoán dương có thể đạt Recall bằng 1 và F1 khá
cao. Baseline Dummy vì thế là đối chứng cần thiết khi đánh giá sự cải thiện.
Các tỷ lệ trên thuộc toàn bộ nguồn, không phải tỷ lệ riêng của train hay test.

\fig[.65]{01_classes.png}{Phân phối nhãn trên tập huấn luyện sau chia nhóm.}{fig:classes}
Hình~\ref{fig:classes} dùng dữ liệu huấn luyện; tỷ lệ dương train trong
artifact là khoảng 53,18\%. Sự khác biệt với tỷ lệ 54,83\% của toàn nguồn
phù hợp với việc chỉ lấy một tập con theo nhóm. Đây là sai khác mô tả,
không phải kiểm định xác nhận mẫu đại diện. Test có 632 dòng dương trên
1.176 dòng và chỉ được dùng đánh giá sau khi khóa lựa chọn mô hình.

Đối chiếu bảng trạng thái gốc cho thấy 7.754 dòng ``Shipping canceled'' vẫn
thuộc lớp 0; 41.592 dòng giao sớm và 32.196 dòng đúng hạn cũng ở lớp này.
Bởi vậy, diễn giải nhãn 0 là ``không được nguồn gán giao trễ'' chính xác hơn
``đã giao thành công đúng hạn''. Không tự loại đơn hủy hoặc tạo lại nhãn vì
điều đó sẽ đổi quần thể và làm kết quả không còn so sánh trực tiếp được với
protocol đã lưu.

\subsection{Chất lượng dữ liệu: thiếu, bất thường và trùng lặp}
\begin{table}[!htbp]\centering
\caption{Giá trị thiếu trên toàn nguồn theo từ điển dữ liệu đã lưu.}\label{tab:missing}
\begin{tabular}{lrr}\toprule
Trường & Số thiếu & Tỷ lệ xấp xỉ\\\midrule
Customer Lname & 8 & 0,004\%\\
Customer Zipcode & 3 & 0,002\%\\
Order Zipcode & 155.679 & 86,24\%\\
Product Description & 180.519 & 100,00\%\\\bottomrule
\end{tabular}

\end{table}
Bảng~\ref{tab:missing} chỉ liệt kê các cột có thiếu. Tất cả đều bị loại khỏi
đầu vào; các đặc trưng được chọn không có missing trong kiểm toán nguồn.
Do vậy, median/mode là cơ chế phòng vệ cho dữ liệu mới, không phải một cải
tiến đã được chứng minh tạo chênh lệch metric trên tập này. Cột mô tả sản
phẩm thiếu hoàn toàn cũng không được coi là dữ liệu văn bản có thể khai thác.

Với trường số, kiểm toán sử dụng quy tắc ngoài khoảng
$[Q_1-1{,}5\,IQR,Q_3+1{,}5\,IQR]$, trong đó $IQR=Q_3-Q_1$. Đây là tiêu
chí đánh dấu, không là kết luận bản ghi sai. Chẳng hạn, lợi nhuận có 18.942
giá trị ngoài ngưỡng IQR nhưng trường này đã bị loại vì ngữ nghĩa và thời
điểm chưa phù hợp. Loại tự động các dòng như vậy có thể vô tình đổi phân
phối nhãn và làm mất đơn có giá trị kinh doanh đáng quan tâm.

Không phát hiện dòng trùng hoàn toàn nên không có căn cứ chạy khử trùng.
Các dòng chung mã đơn phải được giữ cùng nhóm khi chia dữ liệu. Ngoài ra,
CSV được đọc bằng CP1252 trong manifest; giải mã thành công chỉ xác nhận
cú pháp ký tự, chưa bảo đảm mọi địa danh đúng ngữ nghĩa. Không thay byte lỗi
bằng ký tự rỗng để làm chương trình chạy, vì có thể gộp nhầm các danh mục.

\subsection{Thuộc tính quan trọng và ý nghĩa cảnh báo sớm}
Số ngày giao theo kế hoạch mô tả thời hạn cam kết, khác với số ngày giao
thực tế bị cấm sử dụng. Phương thức giao hàng là đặc trưng đáng khảo sát vì
liên quan cách tổ chức dịch vụ. Giá, số lượng, chiết khấu và các nhóm thị
trường có thể phản ánh cấu trúc đơn, nhưng tương quan quan sát được không
chứng minh nguyên nhân của việc giao trễ. Đặc trưng lịch trong luồng B là
một giả thuyết về biến thiên vận hành theo thời gian.

Một hệ thống cảnh báo có ích cần đồng thời xác định đúng thời điểm dữ liệu,
phân biệt nhãn, kiểm soát cảnh báo sai và đánh giá nhóm dễ bị bỏ sót. Vì vậy,
thành công của đề tài không chỉ là tìm một mô hình có điểm lớn nhất mà còn
là chỉ ra trường hợp mô hình không đủ tin cậy. Đây là cơ sở để thiết kế
pipeline và protocol đánh giá trong phần tiếp theo.
